% Generated draft: requires mathematical and source review.
\documentclass[11pt]{article}
\usepackage[margin=25mm]{geometry}
\usepackage{fontspec}
\setmainfont{Latin Modern Roman}
\usepackage{amsmath,amssymb,mathtools,unicode-math}
\setmathfont{Latin Modern Math}
\usepackage{enumitem,needspace,xurl,hyperref}
\setlength{\parindent}{0pt}
\setlength{\parskip}{5pt}
\setlength{\emergencystretch}{3em}
\newcommand{\N}{\mathbb{N}}
\newcommand{\Z}{\mathbb{Z}}
\newcommand{\Q}{\mathbb{Q}}
\newcommand{\R}{\mathbb{R}}
\newcommand{\C}{\mathbb{C}}
\newcommand{\F}{\mathbb{F}}
\newcommand{\A}{\mathbb{A}}
\newcommand{\PP}{\mathbb{P}}
\newcommand{\E}{\mathbb{E}}
\newcommand{\eps}{\varepsilon}
\newcommand{\dd}{\,\mathrm{d}}
\begin{document}
\section*{UnsolvedMath ID: 801050}
\textbf{Source problem code:} catalogue entry
\subsection{Definitions and mathematical statement}
Let \(\{H_\lambda\}_{\lambda\in\Lambda}\) denote the two-dimensional H\'enon family intended in the supplied source. The record does not specify a normal form, the admissible parameter set \(\Lambda\), or an exact convention for an attractor; these remain qualifications of the formulation below.

For a parameter \(\lambda\), an invariant Borel probability measure \(\mu\) satisfies
\[
\mu(H_\lambda^{-1}(B))=\mu(B)
\qquad\text{for every Borel set }B
\]
in the phase space. Following the source's stated description, the SRB property means that the conditional measures along unstable curves are absolutely continuous with respect to length on those curves. More explicitly, for an appropriate measurable partition into local unstable curves \(\gamma\), its conditional measures satisfy
\[
\mu_\gamma\ll m_\gamma
\qquad\text{for almost every partition element }\gamma,
\]
where \(m_\gamma\) denotes induced length. The existence and precise meaning of these unstable curves must be understood in the intended dynamical setting.

\textbf{Question 2.} Give a short proof that there exist a parameter \(\lambda\in\Lambda\), an attractor \(A\) for \(H_\lambda\), and an invariant Borel probability measure \(\mu\) such that
\[
\mu(A)=1,
\qquad
\mu_\gamma\ll m_\gamma
\quad\text{along unstable curves as above}.
\]
Here ``short'' expresses the source's request for a quick proof; the source gives no formal bound on proof length or complexity. The requested advance concerns the proof: the supplied text already reports existence for a positive-measure set of parameters, with difficult and lengthy proofs attributed to references labeled [BC], [Berger], [Ta], and [WY]. Neither the parameter region nor the reference parameter measure is specified in this record.

\textbf{Background and unresolved qualifications.} The source motivates the question by asserting that, for \(q_c(x)=x^2+c\) with \(c\in\R\), an absolutely continuous measure is easy to obtain by choosing a parameter whose post-critical orbit is ``finite but not periodic.'' A finite forward orbit is eventually periodic, so this wording requires clarification, plausibly referring to strict preperiodicity rather than absence of eventual periodicity. This background assertion is retained as supplied context, not as a separately verified theorem. The question occurs in a discussion of stochastic attractors, but the record does not explicitly state positive-Lyapunov-exponent, nontriviality, or dissipation hypotheses. These omissions matter for distinguishing the intended SRB attractor from any interpretation involving attracting periodic orbits.
\subsection{Short English statement}
Find a quick proof that some H\'enon map has an attractor carrying an invariant SRB probability measure, whose conditional measures on unstable curves are absolutely continuous. The supplied source reports that existence for a positive-measure set of parameters is already known through lengthy proofs. The intended parameter range, attractor convention, and nontriviality requirements need clarification.
\subsection{Sources}
\begin{itemize}
\item Local UnsolvedMath registry, permanent ID \(801050\), exact source problem code \texttt{catalogue entry}; title: \emph{Henon Maps: Open Problems --- Question 2}. The supplied registry attributes the problem to Berger and Bedford, with proposed year 2023, and labels it \emph{open}. This label is supplied metadata, not an independently verified current status.
\item Berger and Bedford, \emph{Henon Maps: A List of Open Problems} (2023), Question 2. Supplied source URL: \url{https://arxiv.org/abs/2312.03907}. The record reports access on 2026-10-16 and extraction by \texttt{arxiv\_latex}; these are supplied provenance details, not a claim of independent access or verification.
\item The supplied excerpt attributes the known positive-measure parameter existence result to references labeled [BC], [Berger], [Ta], and [WY]. Their full bibliographic entries and individual links are not included in the record.
\end{itemize}
\end{document}
